\section{Background: Introduction to Piezoelectricity}

%start background here. 

Piezoelectricity is a phenomenon that relates mechanical deformation with charge separation and vice versa.  Piezoelectric compounds are ubiquitous in nature (bone, collagen, proteins, crystals, etc.) and 
industry (zinc oxide, lead zirconate titanate, etc.).  
When piezoelectric materials are compressed or deformed in some
manner, they become polarized and the difference in chemical potentials at opposite surfaces can be measured
as a voltage difference.  
This process is due to either the reorientation of dipoles and their corresponding Weiss domains (local areas of similar dipole density) under applied stress or a local 
change in the environment surrounding the domains.
Conversely, piezoelectric materials are also deformed in response to an applied electric field. This deformation is due
to the field interacting with the molecular electrostatic moments (both static and induced). 
%, and increasing the separation
%between oppositely charged moieties.  
For the purposes of this manuscript, we will be concerned with the linear piezoelectric equations, which pertain to
infinitesimally small applied stresses and electric fields, respectively.   
The theory of piezoelectric response is well-established for crystalline (bulk) systems. 
For more information on piezoelectricity we refer the reader to Ref. \cite{ieeePiezo,Martin1972, Cohen:2008gh}
%We can describe the piezoelectric properties of a crystal via the direct or converse
%piezoelectric effect.  
The {\em direct} piezoelectric effect describes the polarization (electric charge density displacement) $D_{i}$ of the crystal in terms of an applied
stress $T_{jk}$ and electric field $f_j$, %~\eq{eq:direct}.
\begin{equation}
\begin{aligned}
\label{eq:direct}
D_i = d_{ijk}T_{jk}+\epsilon_{ij}f_j \\
\bf{D} = \bf{dT}+\bf{\epsilon f}.
\end{aligned}
\end{equation}
Here and in the following, Einstein sum convention is implied for repeated indices. 
The {\em converse} piezoelectric effect describes the strain $E_{ij}$ 
induced in the material when an external field is applied, %~\eq{eq:convers}.
\begin{equation}
\begin{aligned}
\label{eq:convers}
E_{ij} &= \sigma_{ijkl}T_{kl}+d_{ijk}f_k \\
\bf{E} &= \bf{\sigma T}+\bf{df}. 
\end{aligned}
\end{equation}
%The vector $\bf{D}$ corresponds to the electric charge density displacement. The vector $\bf{f}$ corresponds to an applied electric field.
Here the tensor $\bf{E}$ is known as Green strain tensor. 
The stress tensor, $\bf{T}$, is linked to the strain tensor
by the rank-4 compliance tensor, $\bm{\sigma}$, at zero field. 
% \marginnote{Do we need the compliance tensor?}
For aperiodic systems it is easiest to work with the converse equations by applying an electric field without any
stress. 
%In our calculations, we are interested in applying an electric field (without applying any stress), and approximating the deformation.
%We are therefore solely interested in the piezoelectric tensor $\bf{d}$, and furthermore we focus on the
%converse piezoelectric equation in attempts to approximate a small part of $\bf{d}$--namely \dtt.  The reason we focus on \eq{eq:convers} is because we are interested in
%small systems for which it is easier to apply a field and approximate the deformation of the systems rather than apply a stress and approximate the polarization.  Small systems are 
%ideal when we wish to use electronic structure methods with electron correlation (which is difficult in periodic calculations).  
%
The piezoelectric tensor is then given by 
\begin{equation}
\label{eq:piezoTen}
d_{ijk} = \left(\frac{\partial E_{ij}}{\partial f_k}\right)_{{\bf T} = {\bf 0}}.
\end{equation}
We note that the indices pertaining to strain, ${ij}$, depend on the coordinate system used as a reference frame
for measuring strain (which can be e.g. unit cell parameters or cartesian coordinates, etc.), whereas the index contracted with the
field, $k$, depends on the external (laboratory) reference frame of the applied field.

To understand how to calculate $d_{ijk}$ in practice and to provide some context for the following sections, we 
introduce here selected concepts of strain theory within the framework of continuum mechanics. 
In continuum mechanics, the material is treated as a   
collection of a continuum of infinitesimally small particles. 
Obviously, this is not an appropriate description for a
molecular system. 
Starting from these definitions, we will therefore develop a theory for molecular (i.e., discrete, anisotropic and inhomogeneous) systems. 
%We shall then attempt to cast the problem
%at hand-- namely calculating \dtt-- in terms of strain theory. 
A more complete derivation of the Green strain tensor can be found in the Appendix. 
Our discussion closely follows the presentation from Ref. \cite{contMech}.

Generally, strain analysis is concerned with concepts such as longitudinal strain, shear strain, and volumetric strain.
We will be mostly concerned with longitudinal strain for the discussion that follows.
To this end, consider a body undergoing mechanical deformation (\fig{fig:deform}).
% \marginnote{KEITH: Define vector u.}
\begin{figure}
  \centering
  \includegraphics[width=0.48\textwidth]{deform.pdf}
  \caption{General deformation of a body. Adapted from Ref. \cite{contMech}. The position vectors for the undeformed body, $K_0$, and the deformed body, $K$ are given by ${\bf{X}}$ 
  and ${\bf{x}}$ respectively.  The displacement vector, ${\bf{u}}$ map the position in the undeformed coordinates to a position in the deformed coordinates, and indicate the direction and 
  magnitude in the ``shift'' of an infinitetessimal particle in space as the body deforms.  }
  \label{fig:deform}
\end{figure}
Configurations $K_0$ and $K$ correspond to the undeformed and deformed bodies, respectively. 
Particle positions before and after deformation are denoted by $X$ and $x$, respectively, where we note that these
positions are specified within the same (external) coordinate frame. The vector field ${\bf{u}}({\bf{X}}, t)$ gives the displacement vector from a point the undeformed body 
to a corresponding point in the deformed body.  The displacement vector field will become important later when we define the Green strain tensor.  For molecular systems the 
a discrete displacement vector field exists for each atom in a molecule or a system as said body deforms.  
A central question is how a given material line deforms as the body is deformed from the initial ($K_0$) to the final
($K$) shape. 
For our purposes, the material line can be seen as any line drawn through some arbitrary but continuous path of particles
in the body. For convenience, we show the undeformed material line as a straight line segment $P_{0}Q_{0}$ with length
$s_0$ and the deformed material line as the curve along $PQ$ with length $s$. 
The unit vector along the line segment $P_{0}Q_{0}$ is called $\bld{e}$. 
Based on this picture, one defines the longitudinal strain $\epsilon$, i.e. the strain along the direction of $\bld{e}$,
as the relative change in the length of the material line upon deformation in the limit of infinitesimal line length:
\begin{equation}
\label{eq:longstrain}
\epsilon=\lim_{s_0\to 0}\frac{s-s_0}{s_0}=\frac{ds-ds_0}{ds_0}=\frac{ds}{ds_0}-1
\end{equation}
We note that we used a very similar, though less rigorous definition in our previous papers \cite{Werling0,Werling1} where we used the
percent deformation in bond lengths to approximate the piezoelectric deformation.   
As noted, we are not concerned with volumetric or shear strain here, although the Appendix presents more details about
the Green strain tensor which can be used to completely specify the deformation around a point within a body.

The goal now is to find working equations for calculating the longitudinal strain $\epsilon$ from \eq{eq:longstrain}. 
To this end, one needs to determine the arc length $s$ of the deformed material line. 
Treating $s_0$ as a variable parameter, any point along the undeformed material line $P_{0}K_{0}$ can be
written as $\bld{X}_{0} + s_0 \cdot \bld{e}$, where $\bld{X}_0$ is the starting point of the undeformed material line. The positions
of particles in the deformed bodies can then be defined as functions of the undeformed point as well as a time
coordinate $t$ that determines the progress of the deformation: $\bld{x} = \bld{x}(\bld{X}_{0} + s_0 \cdot\bld{e}, t)$.
The arc length of $s$ as a function of $s_0$ (which we treat as a curve parameter) is then  
\begin{equation}
\label{eq:arclen}
s(s_0) = \int_0^{s_0}\sqrt{\frac{dx_i}{d\bar{s}_0}\frac{dx_i}{d\bar{s}_0}}d\bar{s}_0, 
\end{equation}
which suggests that the arc length derivative with respect to $s_0$ and it's square are
\begin{equation}
\label{eq:dsds0}
  \begin{split}
    \frac{ds(s_0)}{ds_0} &= \sqrt{\frac{dx_i}{ds_0}\frac{dx_i}{ds_0}} \\
     \bigg(\frac{ds}{ds_0}\bigg)^2 &= \frac{dx_i}{d\bar{s}_0}\frac{dx_i}{d\bar{s}_0}
  \end{split}
\end{equation}
The derivative of the deformed arc length with respect to the undeformed length parameter is more easily cast in
terms of a mixed derivative of the deformed coordinates $\bld{x}$ with respect to the undeformed coordinates $\bld{X}$,
\begin{equation} 
\label{eq:dGrad}
  F_{ik} = \frac{\partial x_{i}}{\partial X_{k}},
\end{equation}
to yield 
% Looking back to \fig {fig:deform}, We can let the line element $P_0Q_0$ be called $d{\bf{r}_0}$, with length $s_0$ and unit vector ${\bf{e}}$.
% From this we can then write.
% 
% \begin{equation}
% \label{eq:dXkds0}
%  d{\bf{r}}_0={\bf{e}}\,ds_0=dX_k\,{\bf{e}}_k \quad \Rightarrow \quad dX_k=e_k\,ds_0 \quad \Leftrightarrow \quad e_k=\frac{dX_k}{ds_0}
% \end{equation}
% 
% We will use \eq{eq:dXkds0} later.
% If we consider the vector differential along $s$, tangent to the curve, which we call $d{\bf{r}}$ and has components $dx_i$, we may write
% \begin{equation}
% \label{eq:drds0}
% d{\bf{r}}=\frac{\partial{\bf{r}}}{\partial s_0} \quad \Rightarrow \quad dx_i=\frac{x_i}{\partial{s_0}}ds_0
% \end{equation}
% Now we would like to relate the line elements $d{\bf{r_0}}$ from $K_0$ with $d{\bf{r}}$ in $K$.
% \begin{equation}
% \label{eq:drd0}
% d{\bf{r}}=\frac{\partial{\bf{r}}}{\partial {\bf{r}_0}} \cdot d{\bf{r}_0} \quad  \Leftrightarrow \quad dx_i=\frac{\partial x_i}{\partial X_k}dX_k
% \end{equation}
% 
% We thus have the deformation gradient ${\bf{F}}$.
% \begin{equation}
%  \label{eq:defGrad}
%  {\bf{F}}=\bigtriangledown({\bf{r}})=\frac{\partial{\bf{r}}}{\partial{\bf{r}_0}} \quad \Leftrightarrow  \quad F_{ik}=\frac{\partial x_i}{\partial X_k}
% \end{equation}
% 
% 
% We would like to calculate $(\frac{ds}{ds_0})^2$ from before.  So first we ask what is $\frac{\partial x_i}{\partial s_0}\bigg|_{s_0=0}$?
% \begin{equation}
%  \frac{\partial x_i}{\partial s_0}\bigg|_{s_0=0}=\frac{\partial x_i(X,t)}{\partial X_k}
%  \frac{\partial (X_k+s_0e_k)}{\partial s_0}\bigg|_{s_0=0}\quad \Rightarrow \quad \frac{\partial x_i}{s_0}=
%  \frac{\partial x_i}{\partial X_k}\frac{dX_k}{ds_0}=F_{ik}e_k
% \end{equation}
% 
% In the last step we have used our result from \eq{eq:dXkds0}.  This is analogous to a directional derivative, but instead of using
% the gradient of a scalar quantity, we instead use the deformation gradient which is inherently a matrix.
% 
% We can now write
% \begin{equation}
% \label{eq:dsds02}
%  \bigg(\frac{ds}{ds_0}\bigg)^2=\frac{dx_i}{d\bar{s}_0}\frac{dx_i}{d\bar{s}_0}=
%  (F_{ik}e_k)(F_{il}e_l)= {\bf{e\cdot (F^TF)\cdot e}}={\bf{e\cdot C \cdot e}}
% \end{equation}
% ${\bf{C}}$ is known as the Green deformation tensor. We have come a long way, but recall
% that the displacement of a particle during deformation is given by $\bf{u}(\bf{r}_0, t)$.
% We would like to recast the deformation tensor in terms of the displacement gradient, which
% we call $\bf{H}$
% \begin{equation}
% \label{eq:A}
%  \bf{H} = \frac{\partial {\bf{u}}}{\partial {\bf{r}_0}} \quad \Leftrightarrow \quad H_{ik}=\frac{\partial u_i}{\partial X_k}
% \end{equation}
% Recalling that
% \begin{equation}
%  x_i(X,t) = X_i +u_i(X,t)
% \end{equation}
% we may infer
% \begin{equation}
%  \frac{\partial x_i}{\partial X_k} = \frac{\partial X_i}{\partial X_k}+ \frac{\partial u_i}{\partial X_k}
% \end{equation}
% or
% \begin{equation}
%  F_{ik} = \delta_{ik}+ H_{ik} \quad \Leftrightarrow \quad \bf{F} = \bf{1} + \bf{H}
% \end{equation}
% We recast the Green deformation tensor as follows:
% \begin{equation}
%  {\bf{C}}= {\bf{F^TF}} = ({\bf{1}} + {\bf{H^T}})({\bf{1}}+ {\bf{H}})= {\bf{1}}+ {\bf{H}} + {\bf{H^T}} +{\bf{H^TH}}
% \end{equation}
% The Green strain tensor, ${\bf{E}}$ is defined by:
% \begin{equation}
% \label{eq:H2E}
%  {\bf{E}}= \frac{1}{2}({\bf{H}} + {\bf{H^T}} + {\bf{H^TH}})
% \end{equation}
% which we can represent by elements as:
% \begin{equation}
%  E_{kl}=\frac{1}{2}\bigg(\frac{\partial u_k}{\partial X_l}  + \frac{\partial u_l}{\partial X_k} +
%  \frac{\partial u_i}{\partial X_k}\frac{\partial u_i}{\partial X_l}\bigg)
% \end{equation}
% 
% We have accomplished what we set out to do.  The reason for doing this will be made clear later.
% Our methods for measuring the piezocoefficient are directly tied to understanding strain.
% 
% Hence we can rewrite the Green deformation tensor, ${\bf{C}}$, as follows:
% \begin{equation}
%  {\bf{C}} = {\bf{1}} + 2{\bf{E}}
% \end{equation}

%The derivation is presented in the Appendix, but we may write $\bigg(\frac{ds}{ds_0}\bigg)^2$ as 
% Looking back at \eq{eq:dsds02}, we can rewrite the equation as
\begin{equation}
\label{eq:dsds022}
% \begin{split}
 \left(\frac{ds}{ds_0}\right)^2
 %= e_{i} F_{ik} F_{kj} e_{j}
 %\equiv e_{i} C_{ij} e_{j}
 %\equiv e_{i} (\delta_{ij} + 2 E_{ij}) e_{j},
 = \bld{e} \cdot (\bld{F}^{T} \bld{F}) \cdot \bld{e} \\
%  \end{split}
 \end{equation}
 $(\bld{F}^{T} \bld{F})$ is known as the Green deformation tensor and which we denote by ${\bf{C}}$.  
 Appendix A shows the derivation for the relationship between the Green deformation tensor and the Green Strain tensor, ${\bf{E}}$, but we arrive at  
 \begin{equation}
 \begin{split}
 \left(\frac{ds}{ds_0}\right)^2 = \bld{e} \cdot \bld{C} \cdot \bld{e} \\
 &\equiv \bld{e} \cdot (\bld{1} + 2\bld{E}) \cdot \bld{e}  \\
 &= 1 + 2\bld{e} \cdot \bld{E} \cdot \bld{e},
\end{split}
\end{equation}
using $C = \bld{1} + 2\bld{E}$.
% where we have introduced definitions for the Green deformation tensor 
% \begin{equation}
%   \bld{C} = \bld{F}^{T} \bld{F}
% \end{equation}
% and the Green
% strain tensor 
% \begin{equation}
%   \bld{E} = \frac{1}{2} \left(\bld{C} - \bld{1}\right)
% \end{equation}  
With these definitions, we arrive at a working equation for longitudinal strain along a material line defined by the
unit vector $\bld{e}$ as 
\begin{equation}
\label{eq:longstrain2}
  \epsilon 
  = \frac{ds}{ds_0} - 1
  = \sqrt{ 1 + 2{\bf{e}}\cdot \bf{E} \cdot \bf{e}} - 1
\end{equation}
In the limit of small deformations, the working equation for the longitudinal strain becomes
\begin{equation} 
\label{eq:small}
  \epsilon = \bld{e} \cdot \bld{E} \cdot \bld{e}
\end{equation}
which can be seen from a taylor expansion about the point $\bld{e} \cdot \bld{E} \cdot \bld{e} = 0$
%In the last two equations, ${\bf{C}}$ is the Green Deformation tensor and is ${\bf{E}}$ is the Green Strain Tensor.  In \eq{eq:dsds022} we use
%${\bf{C}}={\bf{1}}+2{\bf{E}}$.
We have not yet defined the elements of the Green strain tensor, ${\bf{E}}$, but they are given by 
\begin{equation}
\label{eq:gst}
 E_{kl}=\frac{1}{2}\bigg(\frac{\partial u_k}{\partial X_l}  + \frac{\partial u_l}{\partial X_k} +
  \frac{\partial u_i}{\partial X_k}\frac{\partial u_i}{\partial X_l}\bigg)
\end{equation}
and is derived fully in the Appendix.  
The Green strain tensor only depends on the derivatives of the displacement vector field with respect to the 
undeformed position vectors. Just as the displacement is a vector field with a unique vector specified for every point in the undeformed body, the Green strain tensor 
is a tensor field that fully specifies the strain around a given point in the deformed body in reference to the undeformed body.  

For convenience later on, we may define the matrix ${\bf{A}}$, and recast the Green strain tensor in terms of ${\bf{A}}$.
\begin{equation}
  \label{eq:A}
  A_{ik} = \frac{\partial u_i}{\partial X_k},
\end{equation}
 and recast the Green strain tensor in terms of ${\bf{A}}$.
 \begin{equation}
 \label{eq:H2E}
  {\bf{E}}= \frac{1}{2} \left( \bld{A} + \bld{A}^{T} + \bld{A}^{T}\bld{A} \right),
\end{equation}
% So far, we have not defined the transformation that describes the deformation of the body from $K_0$ to $K$. It is
% useful to recast \eq{eq:longstrain} in terms of this transformation; the longitudinal strain can then be determined by
% measuring the change in coordinates for a given pair of points. To this
% end, one introduces a vector $\bld{u}$ that defines, for each point along the undeformed material line, how it maps onto
% the corresponding point on the deformed body. 
% Since the deformed material
% line is, in general, curved, the transformation vector length and direction depend, in general, on which point the point along the
% undeformed material line is being transformed. 
% By parameterizing each point on the undeformed material line 
% using $\bld{X}_{0}$ as the starting point, $\bld{e}$ for the direction and $s_0$ as a length parameter, the
% transformation vector can be written as a
% function of these parameters: $\bld{u} = \bld{u}(\bld{X}_{0} + s_{0} \cdot \bld{e})$. 
% Using this definition of the body transformation, the Green strain tensor can be rewritten in terms of the
% transformation vectors $\bld{u}$ (see Appendix A for a derivation):
% \begin{equation}
%  \label{eq:H2E}
%   {\bf{E}}= \frac{1}{2} \left( \bld{A} + \bld{A}^{T} + \bld{A}^{T}\bld{A} \right),
% \end{equation}
% where $\bld{A}$ is the matrix of derivatives between the displacement vectors $\bld{u}$ and the undeformed coordinates
% $\bld{X}$, 
% \begin{equation}
%   \label{eq:A}
%   A_{ik} = \frac{\partial u_i}{\partial X_k}
% \end{equation}
% This formulation is advantageous from a practical perspective, because $\bld{E}$ can now be
% easily approximated numerically by measuring displacements. 
% The Green strain tensor could also be used to define shear and volumetric strain in very similar fashions, but we focus
% here on the longitudinal strain.  
We now wish to show how we can approximate \dtt in a manner equivalent to our previous papers \cite{Werling0,Werling1}.
For our situation, we take the undeformed body, $K_0$ to be the optimized geometry at zero field.  
The only source of deformation for our system will be the applied electric field ${\bld{f}}$, which implies that the Green strain tensor
is also a function of ${\bld{f}}$, i.e. $\bld{E} = \bld{E}(\bld{f})$.  We can therefore rewrite \eq{eq:longstrain2}.
\begin{equation} 
  \epsilon = \sqrt{1 + 2 \, \bld{e} \cdot \bld{E}(\bld{f}) \cdot \bld{e}} - 1
                \approx \bld{e} \cdot \bld{E}(\bld{f}) \cdot \bld{e}
\end{equation}
If we take the derivative with respect to $\bld{f}$ we obtain the following.
\begin{equation}
\label{eq:longstrainDeriv}
\frac{\partial \epsilon ({\bf{f}})}{\partial {\bf{f}}} = \frac{{\bf{e}}\cdot \frac{\partial {\bf{E}}({\bf{f}})}{\partial {\bf{f}}} \cdot {\bf{e}}}{\sqrt{ 1 + 2{\bf{e}}\cdot \bf{E(f)} \cdot \bf{e}}}
\end{equation}
Because we choose $K_0$ as our system at zero field, $\bld{E}(\bld{f}=\bld{0})=\bld{0}$ ($\bld{0}$ on the right hand side is a matrix not a vector)--ie. the displacement derivatives contained
in $\bld{A}$ are all zero since the $\bld{u}=\bld{0}$ everywhere. 
If we evaluate \eq{eq:longstrainDeriv} at $\bld{f}=\bld{0}$ the derivative reduces to the following equation.    
\begin{equation}
\label{eq:longstrainDeriv2}
\frac{\partial \epsilon ({\bf{0}})}{\partial {\bf{f}}} = {\bf{e}}\cdot \frac{\partial {\bf{E}}({\bf{0}})}{\partial {\bf{f}}} \cdot {\bf{e}}
\end{equation}
\eq{eq:longstrainDeriv2} not surprisingly matches the field derivative of the longitudinal strain at small defomations (\eq{eq:small}). 
The derivative of the Green strain tensor $\bld{E}$ with respect to the field, $\bld{f}$ is just the piezoelectric tensor, $\bld{d}$ (\eq{eq:piezoTen}). 
Thus the derivative of the longitudinal strain at zero field reduces to the following:
\begin{equation}
\label{eq:longstrainDeriv3}
\frac{\partial \epsilon ({\bf{0}})}{\partial {\bf{f}}} = {\bf{e}}\cdot \bld{d}({\bf{0}}) \cdot {\bf{e}}
\end{equation}
In the sections to come, we will reduce \eq{eq:longstrainDeriv3} to a simpler form, devise a method for calculating piezoelectric coefficients for molecules,
and introduce the piezoelectric matrix, $\bld{P}$.  \eq{eq:longstrainDeriv3} can be used to approximate \dtt in the same manner as our previous papers \cite{Werling0,Werling1}, but is a generalization, which
we will find useful later.  Furthermore, it connects the work we have done thus far to strain theory, which we hope serves as a tool for future work in this area.  


% We can reduce \eq{longstrainDeriv2} further by using the definition for the Green strain tensor $\bld{E}$ (\eq{eq:H2E}). 
% We first rewrite \eq{eq:longstrainDeriv2} elementwise (the derivative is a vector quantity with components of field) using the definition of $\bld{A}$ (\eq{eq:A}).
% \begin{equation}
% \begin{split}
% \label{eq:element}
% \frac{\partial \epsilon ({\bf{0}})}{\partial f_l}&= \frac{1}{2}e_i \bigg(\frac{\partial^2 u_i(\bld{0})}{\partial X_k\partial f_l} + \frac{\partial^2 u_k(\bld{0})}{\partial X_i \partial f_l} + \\
%  &\frac{\partial^2 u_j(\bld{0})}{\partial X_i \partial f_l}\frac{\partial u_j(\bld{0})}{\partial X_k} + \frac{\partial u_j(\bld{0})}{\partial X_i}\frac{\partial^2 u_j(\bld{0})}{\partial X_k \partial f_l}\bigg)e_k
%  \end{split}
% \end{equation}
% The last two terms in \eq{eq:element} result from the product rule for $\bld{A}^T\bld{A}$ from \eq{eq:H2E}.  Both terms vanish when evaluated at zero field.  This is because 
% $\frac{\partial u_j(\bld{0})}{\partial X_i}=A_{ji}(\bld{0})=0$ (every element is zero) if we take the zero field body as $K_0$ and $K$ (which is equivalent to evaluating $\bld{A}$ at zero field. 
% If this is the case, the displacement vector for every point in the body is constant (and actually $\bld{0}$), and hence the derivative in any direction ($X_i$) vanishes (see Appendix
% for more information on strain theory).  We can therefore rewrite \eq{eq:element}, ignoring the last two terms.  
% \begin{equation}
% \label{eq:element2}
% \frac{\partial \epsilon ({\bf{0}})}{\partial f_l}= \frac{1}{2}e_i \bigg(\frac{\partial^2 u_i(\bld{0})}{\partial X_k\partial f_l} + \frac{\partial^2 u_k(\bld{0})}{\partial X_i \partial f_l} \bigg)e_k
% \end{equation}
% We can rewrite \eq{eq:element2} in terms of matrix derivatives for simplicity.  
% \begin{equation}
% \begin{split}
% \label{eq:nonelement}
% \frac{\partial \epsilon ({\bf{0}})}{\partial \bld{f}}&= \frac{1}{2}\bld{e} \bigg(\frac{\partial \bld{A}^T(\bld{0})}{\partial \bld{f}} + \frac{\partial \bld{A}(\bld{0})}{\partial \bld{f}} \bigg)\bld{e} \\
% \frac{\partial \epsilon ({\bf{0}})}{\partial \bld{f}}&= \bld{e}\frac{\partial \bld{A}(\bld{0})}{\partial \bld{f}}\bld{e}
% \end{split}
% \end{equation}
% Because both matrices are contracted by $\bld{e}$ on both sides, and one is the transpose of the other in these two dimensions, both terms in the first line were equivalent and reduced to twice
% the value of one term, which cancelled the prefactor of $\frac{1}{2}$ in the final line.
% $\frac{\partial \bld{A}(\bld{0})}{\partial \bld{f}}$ is a 3 dimensional tensor, and it is equivalent to the derivative of the Green strain tensor evaluated at zero field (ie. the piezoelectric tensor evaluated
% at zero field).  
% We will work with \eq{eq:nonelement} in the sections to come.  










% To arrive at working equations for studying piezoelectric systems, we recall that the piezoelectric tensor
% $\bld{d}$ is defined as the derivative of longitudinal strain (along a given material line) with respect to electric field coordinates. In general,
% the mechanical (geometric) response of the material depends on the field, i.e. $\bld{E} = \bld{E}(\bld{f})$.   
% Defining a set of unit vectors $\{\bld{e}_{i}\}$ that spans the coordinate system in which to consider the piezoelectric
% response, the strain along these axes becomes
% \begin{equation} 
%   \epsilon_{ij} = \sqrt{1 + 2 \, \bld{e}_{i} \cdot \bld{E}(\bld{f}) \cdot \bld{e}_{j}} - 1
%                 \approx \bld{e}_{i} \cdot \bld{E}(\bld{f}) \cdot \bld{e}_{j}
% \end{equation}
% and the corresponding piezoelectric tensor becomes 
% \begin{equation} 
%   d_{ijk} 
%   = \left( \frac{\partial \epsilon_{ij}}{\partial f_{k}} \right)_{\bld{T} = \bld{0}}
%   = \bld{e}_{i} \cdot \frac{\partial E}{\partial f_{k}}(\bld{f}) \cdot \bld{e}_{j}
% \end{equation}
% From this formulation it becomes clear that the piezoelectric response can be readily evaluated by determining the field
% derivatives of the Green strain tensor. 

%Since we are interested in quantifying the piezoelectric response along different coordinate axes, 
%we must concern ourselves with the derivative with respect to the field applied to the system.  For our purposes the distortion represented by the strain tensor solely depends on the 
%applied field.  We can therefore write the following:
%\begin{equation}
%\label{eq:longstrain2}
%\epsilon=\frac{ds}{ds_0}-1= \sqrt{ 1 + 2{\bf{e}}\cdot {\bf{E}}({\bf{f}}) \cdot \bf{e}} -1
%\end{equation}
%where ${\bf{f}}$ is the field vector applied to the system.  If we take the derivative of $\epsilon$ with respect to the field vector we obtain
%\begin{equation}
%\label{eq:longstrainDeriv}
%\frac{\partial \epsilon ({\bf{f}})}{\partial {\bf{f}}} = \frac{{\bf{e}}\cdot \frac{\partial {\bf{E}}({\bf{f}})}{\partial {\bf{f}}} \cdot {\bf{e}}}{\sqrt{ 1 + 2{\bf{e}}\cdot \bf{E(f)} \cdot \bf{e}}}
%\end{equation}
%We are interested in the piezoelectric properties of the molecule at zero field.  Without writing the form of ${\bf{E}}({\bf{f}})$ explicitly (although this is easy if we are concerned only with the 
%linear piezoelectric effect), we may assert that ${\bf{E}}({\bf{0}})=0$.  We therefore evaluate ~\eq{longstrainDeriv} for zero field.
%\begin{equation}
%\label{eq:longstrainDeriv2}
%\frac{\partial \epsilon ({\bf{0}})}{\partial {\bf{f}}} = {\bf{e}}\cdot \frac{\partial {\bf{E}}({\bf{0}})}{\partial {\bf{f}}} \cdot {\bf{e}}
%\end{equation}
%~\eq{eq:longstrainDeriv2} is what we must calculate to approximate \dtt.  
%%ends here


% \subsection{Continuum Mechanics (Stress and Strain)}
% In the sections to follow, we will present the linear piezoelectric equation.  Before this, we will
% give a brief overview of some areas of continuum mechanics (most notably stress and strain) which we will find useful later.
% The discussions here will be adapted from ~\cite{contMech}.
% 
% Continuum mechanics is rooted in the continuum hypothesis.
% The continuum hypothesis in summary states that matter is infinitely divisible and no matter how small, a volume element of the matter
% will always contain matter.  Also, for a given type of matter, the matter will maintain continuity throughout deformation and movement ~\cite{contMech}.
% We know that these statements are not rigorously true, since matter is made up of discrete objects like atoms and molecules, and there
% can indeed be a great deal of space (in a relative sense) between atoms and molecules in a material.  However, in general the great number of
% of atoms and molecules which exist in a relevant size (macroscopic world) of some solid or liquid body, would lead to the conclusion that
% this hypothesis is ``good enough'' to work with.  And actually, for our purposes, the violation of the continuuity of a crystal body will
% make our work with the linear piezotensor much easier.
% 
% \begin{figure}
% 
%   \centering
%   \includegraphics[width=\textwidth]{stressUnitCube.png}
%   \caption{Unit of continuum material ~\cite{contMech}.}
%   \label{fig:unit}
% \end{figure}
% 
% We will start our review with an infinitesimal unit of a some continuum material (see \fig {fig:unit}). For our purposes we will consider
% the stress vectors ($\bf{t}$)acting on the unit and the strain which results.  Stress has units of pressure and can be imagined as a pressure vector
% acting on or across a surface. We notice that the author uses ${\bf{t}}$ to denote the individual stress vectors acting on the 3 ``front'' surfaces of the unit cell.
% We also notice that there is a unit vector normal ($\bf{e}_i$) to each surface which make up our basis.  In \fig {fig:unit}, we see that the individual stress vectors ${\bf{t_i}}$ can be
% broken up into
% component stresses (in our basis) ${\bf{T_{ji}}}$ along the face of each front surface.
% 
% At this point we should mention that there are stress vectors that act on the faces of the opposing sides of our ``cube''
% of material.  However, they will be equal and opposite with respect to the opposing surface and we need not consider them to
% fully describe the stress in the material~\cite{contMech}.
% 
% We can furthermore organize this information into a matrix ${\bf{T}}$.  The diagonal of the matrix ${\bf{T_{ii}}}$ consists of the normal stresses.
% These stresses point ``away from (or into) the surface'' and in the case of a rectrangular system are normal to the surface.
% The off diagonal elements, ${\bf{T_{ij}}}, i\neq j$ consist of the shear stresses which generally act along the surface of the piece of the continuum
% material.  Stress, in 3 dimensions, has units of pressure.
% 
% It turns out that this matrix contains all the information we need about the stress acting on a particle within a body.  The Cauchy stress-theorem, which
% we will not prove, states that the stress vector acting on a surface through our unit cube is uniquely determined by our stress tensor of the particle
% and the normal of the surface.  In other words, if we were to ``cut'' through our differential element, creating some exposed face with unit normal ${\bf{n}}$, we could
% calculate the stress vector ${\bf{t}}$ acting on the surface via ${\bf{t}}= {\bf{T}} \dot {\bf{n}}$~\cite{contMech}. This is all we need to know about the stress tensor for the time
% being and we can move on to the strain tensor.
% 
% Deformation analysis generally concerns the local deformation around a point in a continuum material.  We are interested in the strain tensor ${\bf{S}}$, which we will use to understand
% piezoelectric effect in the following section.  Deformation analysis is a rich subject with many topics, but we shall attempt to stick to what is necessary for the present work
% and following section.
% 
% Our discussion will follow the line of reasoning presented in~\cite{contMech}.
% Generally, strain analysis is concerned with concepts like longitudinal strain, shear strain, and volumetric strain.
% These have simple definitions, but first let us consider a body undergoing deformation (\fig {fig:deform}).
% 
% \begin{figure}
%   \centering
% 
%   \includegraphics[width=\textwidth]{deform.png}
%   \caption{General deformation of a body.~\cite{contMech}}
%   \label{fig:deform}
% \end{figure}
% 
% Notice that we have configuration $K_0$ and $K$ corresponding to the undeformed and deformed bodies respectively, and particle positions are given by $X$ and $x$ respectively
% for each configuration.  We can define a material line, which in our figure is given by the line segment $P_0Q_0$ with length $s_0$ in $K_0$, and ask what happens to the material line as the body deforms.
% For our purposes the material line can be seen as a line drawn through some arbitrary group of particles in the body (here it is a straight line for convenience), and we are
% asking what the line will look like if we draw it through the same particles after a deformation, which is given by the curve along $PQ$ with length $s$ in $K$. The deformed material line is,
% in general, curved.
% Looking at this picture we can define the longitudinal strain in the direction of $\bf{e}$, which we will call $\epsilon$ as follows:
% \begin{equation}
% \label{eq:longstrain}
% \epsilon=\lim_{s_0\to 0}\frac{s-s_0}{s_0}=\frac{ds-ds_0}{ds_0}=\frac{ds}{ds_0}-1
% \end{equation}
% This can be viewed as the change in length of the material line per length of the the undeformed line segment $P_0Q_0$.
% We actually use an approximation to this value later when we attempt to calculate piezocoefficients for our systems.
% 
% We will not concern ourselves with volumetric and shear strain for now.
% We want to arrive at the Green strain tensor which can be used to completely specify the deformation around a point in our body.
% First we will allow $s_0$ to be a curve parameter which we may vary for the material lines in both $K_0$ and $K$. We can therefore specify
% points along each material line from $P_0$ to $Q_0$ as $X_i+s_0e_i$, which we trace out as $s_0$ increases. We can therefore express the positions
% of particles in the deformed bodies as a function of this line segment and time--$x_i({\bf{X}}+s_0{\bf{e}},t)$.
% 
% Evaluating the arc length of $s$ as a function of $s_0$ which we treat as a curve parameter, we find that
% 
% \begin{equation}
% \label{eq:arclen}
% s(s_0) = \int_0^{s_0}\sqrt{\frac{dx_i}{d\bar{s}_0}\frac{dx_i}{d\bar{s}_0}}d\bar{s}_0
% \end{equation}
% 
% This indicates to us that
% 
% \begin{align}
% \label{eq:dsds0}
% \frac{ds(s_0)}{ds_0} &= \sqrt{\frac{dx_i}{ds_0}\frac{dx_i}{ds_0}} \\
% ds^2 &= \frac{dx_i}{ds_0}\frac{dx_i}{ds_0}ds_0^2
% \end{align}
% 
% Looking back to \fig {fig:deform}, We can let the line element $P_0Q_0$ be called $d{\bf{r}_0}$, with length $s_0$ and unit vector ${\bf{e}}$.
% From this we can then write.
% 
% \begin{equation}
% \label{eq:dXkds0}
%  d{\bf{r}}_0={\bf{e}}\,ds_0=dX_k\,{\bf{e}}_k \quad \Rightarrow \quad dX_k=e_k\,ds_0 \quad \Leftrightarrow \quad e_k=\frac{dX_k}{ds_0}
% \end{equation}
% 
% We will use \eq{eq:dXkds0} later.
% If we consider the vector differential along $s$, tangent to the curve, which we call $d{\bf{r}}$ and has components $dx_i$, we may write
% \begin{equation}
% \label{eq:drds0}
% d{\bf{r}}=\frac{\partial{\bf{r}}}{\partial s_0} \quad \Rightarrow \quad dx_i=\frac{x_i}{\partial{s_0}}ds_0
% \end{equation}
% Now we would like to relate the line elements $d{\bf{r_0}}$ from $K_0$ with $d{\bf{r}}$ in $K$.
% \begin{equation}
% \label{eq:drd0}
% d{\bf{r}}=\frac{\partial{\bf{r}}}{\partial {\bf{r}_0}} \cdot d{\bf{r}_0} \quad  \Leftrightarrow \quad dx_i=\frac{\partial x_i}{\partial X_k}dX_k
% \end{equation}
% 
% We thus have the deformation gradient ${\bf{F}}$.
% \begin{equation}
%  \label{eq:defGrad}
%  {\bf{F}}=\bigtriangledown({\bf{r}})=\frac{\partial{\bf{r}}}{\partial{\bf{r}_0}} \quad \Leftrightarrow  \quad F_{ik}=\frac{\partial x_i}{\partial X_k}
% \end{equation}
% 
% 
% We would like to calculate $(\frac{ds}{ds_0})^2$ from before.  So first we ask what is $\frac{\partial x_i}{\partial s_0}\bigg|_{s_0=0}$?
% \begin{equation}
%  \frac{\partial x_i}{\partial s_0}\bigg|_{s_0=0}=\frac{\partial x_i(X,t)}{\partial X_k}
%  \frac{\partial (X_k+s_0e_k)}{\partial s_0}\bigg|_{s_0=0}\quad \Rightarrow \quad \frac{\partial x_i}{s_0}=
%  \frac{\partial x_i}{\partial X_k}\frac{dX_k}{ds_0}=F_{ik}e_k
% \end{equation}
% 
% In the last step we have used our result from \eq{eq:dXkds0}.  This is analogous to a directional derivative, but instead of using
% the gradient of a scalar quantity, we instead use the deformation gradient which is inherently a matrix.
% 
% We can now write
% \begin{equation}
% \label{eq:dsds02}
%  \bigg(\frac{ds}{ds_0}\bigg)^2=\frac{dx_i}{d\bar{s}_0}\frac{dx_i}{d\bar{s}_0}=
%  (F_{ik}e_k)(F_{il}e_l)= {\bf{e\cdot (F^TF)\cdot e}}={\bf{e\cdot C \cdot e}}
% \end{equation}
% ${\bf{C}}$ is known as the Green deformation tensor. We have come a long way, but recall
% that the displacement of a particle during deformation is given by $\bf{u}(\bf{r}_0, t)$.
% We would like to recast the deformation tensor in terms of the displacement gradient, which
% we call $\bf{H}$
% \begin{equation}
% \label{eq:A}
%  \bf{H} = \frac{\partial {\bf{u}}}{\partial {\bf{r}_0}} \quad \Leftrightarrow \quad H_{ik}=\frac{\partial u_i}{\partial X_k}
% \end{equation}
% Recalling that
% \begin{equation}
%  x_i(X,t) = X_i +u_i(X,t)
% \end{equation}
% we may infer
% \begin{equation}
%  \frac{\partial x_i}{\partial X_k} = \frac{\partial X_i}{\partial X_k}+ \frac{\partial u_i}{\partial X_k}
% \end{equation}
% or
% \begin{equation}
%  F_{ik} = \delta_{ik}+ H_{ik} \quad \Leftrightarrow \quad \bf{F} = \bf{1} + \bf{H}
% \end{equation}
% We recast the Green deformation tensor as follows:
% \begin{equation}
%  {\bf{C}}= {\bf{F^TF}} = ({\bf{1}} + {\bf{H^T}})({\bf{1}}+ {\bf{H}})= {\bf{1}}+ {\bf{H}} + {\bf{H^T}} +{\bf{H^TH}}
% \end{equation}
% The Green strain tensor, ${\bf{E}}$ is defined by:
% \begin{equation}
% \label{eq:H2E}
%  {\bf{E}}= \frac{1}{2}({\bf{H}} + {\bf{H^T}} + {\bf{H^TH}})
% \end{equation}
% which we can represent by elements as:
% \begin{equation}
%  E_{kl}=\frac{1}{2}\bigg(\frac{\partial u_k}{\partial X_l}  + \frac{\partial u_l}{\partial X_k} +
%  \frac{\partial u_i}{\partial X_k}\frac{\partial u_i}{\partial X_l}\bigg)
% \end{equation}
% 
% We have accomplished what we set out to do.  The reason for doing this will be made clear later.
% Our methods for measuring the piezocoefficient are directly tied to understanding strain.
% 
% 
% Hence we can rewrite the Green deformation tensor, ${\bf{C}}$, as follows:
% \begin{equation}
%  {\bf{C}} = {\bf{1}} + 2{\bf{E}}
% \end{equation}
% Looking back at \eq{eq:dsds02}, we can rewrite the equation as
% \begin{equation}
% \label{eq:dsds022}
%  \bigg(\frac{ds}{ds_0}\bigg)^2={\bf{e\cdot C \cdot e}}={{\bf{e}} \cdot ({\bf{1}} + 2{\bf{E}})\cdot {\bf{e}}} =
%  1 + 2{\bf{e}\cdot \bf{E} \cdot \bf{e}}
% \end{equation}
% Hence, recalling \eq{eq:longstrain}, we can also rewrite our longitudinal strain from before in terms of the Green strain tensor.
% \begin{equation}
% \label{eq:longstrain2}
% \epsilon=\frac{ds}{ds_0}-1= \sqrt{ 1 + 2{\bf{e}}\cdot \bf{E} \cdot \bf{e}} -1
% \end{equation}
% 
% We can also formulate ideas about shear and volumetric strain in terms of the Green strain tensor.  But for
% now we just understand that we have developed our ideas about strain enough to utilize them for studying
% our piezoelectric systems.
% 
% One more brief topic of discussion is principal strain.  Without going into too much detail, we can understand
% principal strain as the material lines which maximize or minimize the value of longitudinal strain.  We can quickly demonstrate this.
% The equation for longitudinal strain is given by \eq{eq:longstrain2} in terms of the Green strain tensor.  Looking at the equation carefully,
% reveals that finding the extrema for longitudinal strain is the same as finding the extrema for $\bf{e}\cdot \bf{E} \cdot \bf{e}$, since the function
% is monotonic in this term.  We therefore add the Lagrange multiplier term $\lambda (\bf{e}\cdot\bf{e}-1)$
% to the strain tensor term, where $\lambda$ is the Lagrange multiplier.  This ensures that the direction vector $\bf{e}$ is normalized.
% \begin{equation}
% \label{eq:strainOpt}
%  {\bf{e}}\cdot {\bf{E}} \cdot {\bf{e}}-\lambda ({\bf{e}}\cdot{\bf{e}}-1)
% \end{equation}
% 
% Since we want to find the extremum of this expression we take the gradient with respect to $\bf{e}$ and set the value equal to the zero vector.
% This results in the following equation:
% \begin{align}
%  \bf{E}\cdot\bf{e} -\lambda \bf{e} &= \bf{0} \\
%   \bf{E}\cdot\bf{e} &= \lambda \bf{e}
% \end{align}
% This is a typical eigenvalue problem.  We note from the above work that the strain
% tensor is symmetric and therefore the eigenvectors are orthogonal and the eigenvalues real.  By diagonalizing the strain tensor,
% we find the directions of principal strain and the maximum and minimum longitudinal strain.  If we call our largest eigenvalue $\tilde{\epsilon_1}$ for example, the resulting
% principal strain would be given by
% \begin{equation}
%  \epsilon_1 = \sqrt{1+2\tilde{\epsilon_1}} -1
% \end{equation}
% after substituting into \eq{eq:longstrain}.
% This will be important later when we try to derive the formula to determine the largest possible \dtt.




% \subsection{The Linear Piezoelectric Equation}
% We now have the necessary information to discuss the linear piezoelectric equation. The equations in this section come from ~\cite{ieeePiezo}The piezoelectric effect
% can be described via the direct or converse piezoelectric effect.  We can relate the strain tensor to the piezoelectric tensor in the following ways.
% 
% For the direct piezoelectric effect we have
% \begin{align}
% \label{eq:direct}
% D_i &= d_{ijk}T_{jk}+\epsilon_{ij}E_j \\
% \bf{D} &= \bf{dT}+\bf{\epsilon E}
% \end{align}
% and for the converse piezoelectric effect we have
% \begin{align}
% \label{eq:convers}
% S_{ij} &= \sigma_{ijkl}T_{kl}+d_{ijk}E_k \\
% \bf{S} &= \bf{\sigma T}+\bf{dE}
% \end{align}
% 
% 
% The vector $\bf{D}$ corresponds to the electric charge density displacement. The vector $\bf{E}$ corresponds to an applied electric field.
% The tensor $\bf{S}$ is strain (the Green strain tensor), which we know to be of rank 2.  The stress tensor, $\bf{T}$, which is of rank 2, is linked to the strain tensor
% by the compliance tensor, $\bf{\sigma}$, rank 4. In our calculations, we are interested in applying a field and approximating the deformation.
% We are therefore solely interested in the piezoelectric tensor $\bf{d}$, and furthermore we focus on the
% converse piezoelectric equation in attempts to approximate a small part of $\bf{d}$--namely \dtt.  The reason we focus on \eq{eq:convers} is because we are interested in
% small systems for which it is easier to apply a field and approximate the deformation of the systems rather than apply a stress and approximate the polarization.
% 
% 
% Based on the equation for the converse piezoelectric effect, we can define the piezoelectric tensor as:
% \begin{equation}
% \label{eq:piezoTen}
% \frac{\partial S_{ij}}{\partial E_k}=d_{ijk}
% \end{equation}
% It is worthwhile to mention that we have an alternative definition for the linear piezoelectric tensor based on the direct equation.
% \begin{equation}
% \label{eq:piezoTen2}
% \frac{\partial D_{i}}{\partial T_{jk}}=d_{ijk}
% \end{equation}
% This definition is most likely more suitable for calculations for periodic crystals.
% 
% We should point out that the indexes pertaining to strain, ${ij}$, depend on the coordinate system we wish to use as a reference frame
% for our normal and shear stresses, but the index contracted with the field, $k$, is dependent on the reference frame for the applied field.


